\chapter{方法与实验}
\section{数学表达与交叉引用}
变量一般采用斜体，单位符号采用正体。公式按章编号，例如式\eqref{eq:sample}：
\begin{equation}
  E = mc^2 .\label{eq:sample}
\end{equation}
此处仅演示公式排版，不构成论文研究内容。

\section{图与表}
图表应先在正文中提及，再随文出现。图\ref{fig:sample}展示示意图的占位方式，表\ref{tab:sample}展示三线表格式。
\begin{figure}[htbp]
  \centering
  \fbox{\parbox[c][22mm][c]{0.75\linewidth}{\centering 在此替换为研究图像}}
  \caption{示意图占位示例}\label{fig:sample}
\end{figure}
\begin{table}[htbp]
  \centering
  \caption{实验结果排版示例}\label{tab:sample}
  \begin{tabular}{lcc}
    \toprule
    方法 & 指标一 & 指标二\\
    \midrule
    方法甲 & 0.00 & 0.00\\
    方法乙 & 0.00 & 0.00\\
    \bottomrule
  \end{tabular}
\end{table}

\section{实验分析}
填写实验设置、评价方法和结果分析。不得将示例数值作为真实实验结论使用。
